% =========================================================================
% myIU thesis -- Nguyen Tan Phat, ITITIU21354
% "myIU: Design and Implementation of an Integrated University Student
% Portal with Federated Identity Authentication and Dual-Interface
% Management Architecture"
% International University - VNU-HCM, School of Computer Science and
% Engineering.
%
% EDITOR'S NOTES (compile-time comments only; nothing in this block prints):
%   - Format/preamble follows the same convention established for the
%     Huynh Ngoc Anh Thu (MediMaster) thesis rewrite in this session:
%     documentclass[12pt,a4paper]{report}, plain left=1.5in/right=1in/
%     top=1in/bottom=1in margins, default serif font, indentfirst,
%     fancyhdr centered folio, \section*{} for Acknowledgments/Abstract,
%     real auto-numbered \chapter/\section/\subsection and \caption'd
%     figures/tables so numbering is 100% automatic. NONE of Anh Thu's
%     content was reused -- only the presentation style/structure.
%   - Content is drawn from this repository's own documentation, written
%     and verified against the actual codebase before this file existed:
%     docs/thesis-outline.md (chapter structure), docs/thesis.html
%     (detailed prose + verified code excerpts), docs/system-design.html
%     (Mermaid architecture/ERD/class/sequence diagrams), docs/API_ROUTES.md
%     (endpoint reference), docs/evaluation-report.md (real test-run and
%     security-finding numbers, generated by actually running the test
%     suites and linters, not from memory).
%   - Diagrams: NONE are hand-drawn/fabricated by this rewrite. The
%     source docs already contain 13 real Mermaid diagrams in
%     docs/system-design.html; every \includegraphics{} call below is a
%     placeholder pointing at Images/diagram-X-X-slug.png with a LaTeX
%     comment directly above it stating exactly which numbered diagram
%     in system-design.html to screenshot (Mermaid renders in-browser --
%     open the file, screenshot the rendered SVG, save under that exact
%     filename into an Images/ folder next to this .tex file, and upload
%     that whole folder to Overleaf). Entity/relationship attribute
%     tables are additionally kept alongside the ERD screenshots
%     (matching how Anh Thu's own thesis already relies heavily on
%     entity tables alongside its figures).
%   - Acknowledgments and the exact defense date are personal information
%     this rewrite cannot know; placeholders are marked clearly -- please
%     edit them before submitting.
%   - Chapter 5's security/testing findings are the same ones already
%     verified in docs/evaluation-report.md (real command output, not
%     estimated), reproduced here with the same honesty about what has
%     and has not been done.
% =========================================================================

\documentclass[12pt,a4paper]{report}

\usepackage{amsmath,amssymb}
\usepackage{iftex}
\ifPDFTeX
  \usepackage[T1]{fontenc}
  \usepackage[utf8]{inputenc}
  \usepackage{textcomp}
\else
  \usepackage{unicode-math}
  \defaultfontfeatures{Scale=MatchLowercase}
\fi

\usepackage{graphicx}
\graphicspath{{Images/}}
\usepackage{titlesec}
\usepackage{fancyhdr}
\usepackage{tocloft}
\usepackage{setspace}
\usepackage{geometry}
\usepackage{hyperref}
\usepackage{indentfirst}

\geometry{left=1.5in, right=1in, top=1in, bottom=1in}
\pagestyle{fancy}
\fancyhf{}
\fancyfoot[C]{\thepage}

\IfFileExists{upquote.sty}{\usepackage{upquote}}{}
\IfFileExists{microtype.sty}{\usepackage{microtype}}{}

\usepackage{xcolor}
\usepackage{longtable,booktabs,array}
\usepackage{calc}
\usepackage{etoolbox}
\makeatletter
\patchcmd\longtable{\par}{\if@noskipsec\mbox{}\fi\par}{}{}
\makeatother
\IfFileExists{footnotehyper.sty}{\usepackage{footnotehyper}}{\usepackage{footnote}}
\makesavenoteenv{longtable}

\makeatletter
\def\maxwidth{\ifdim\Gin@nat@width>\linewidth\linewidth\else\Gin@nat@width\fi}
\def\maxheight{\ifdim\Gin@nat@height>\textheight\textheight\else\Gin@nat@height\fi}
\makeatother
\setkeys{Gin}{width=\maxwidth,height=\maxheight,keepaspectratio}
\makeatletter
\def\fps@figure{htbp}
\makeatother

\usepackage{caption}
\captionsetup{font=small,labelfont=bf,tableposition=top}

\usepackage{algorithm}
\usepackage{algpseudocode}

\setlength{\emergencystretch}{3em}
\providecommand{\tightlist}{\setlength{\itemsep}{0pt}\setlength{\parskip}{0pt}}

\ifLuaTeX
  \usepackage{selnolig}
\fi
\IfFileExists{xurl.sty}{\usepackage{xurl}}{}
\urlstyle{same}
\hypersetup{hidelinks,pdfcreator={LaTeX}}

\author{}
\date{}

\begin{document}

\pagenumbering{roman}

% ------------------------------------------------------------------ %
% COVER PAGE
% ------------------------------------------------------------------ %
\begin{titlepage}
    \centering
    \vspace*{0.2in}
    \textbf{VIETNAM NATIONAL UNIVERSITY OF HO CHI MINH CITY}\\
    \textbf{THE INTERNATIONAL UNIVERSITY}\\
    \textbf{SCHOOL OF COMPUTER SCIENCE AND ENGINEERING}\\[0.4in]
    \textbf{myIU: DESIGN AND IMPLEMENTATION OF AN INTEGRATED UNIVERSITY STUDENT PORTAL WITH FEDERATED IDENTITY AUTHENTICATION AND DUAL-INTERFACE MANAGEMENT ARCHITECTURE}\\[0.2in]
    By\\[0.2in]
    \textbf{Nguyen Tan Phat - ITITIU21354}
    \vfill
    \textit{A thesis submitted to the School of Computer Science and Engineering}\\
    \textit{in partial fulfillment of the requirements for the degree of}\\
    \textit{Engineer of Information Technology}
    \vfill
    Ho Chi Minh City, Vietnam\\
    2025
\end{titlepage}

% ------------------------------------------------------------------ %
% APPROVAL PAGE
% ------------------------------------------------------------------ %
\newpage
\thispagestyle{empty}
\begin{center}
\textbf{myIU: DESIGN AND IMPLEMENTATION OF AN INTEGRATED UNIVERSITY STUDENT PORTAL WITH FEDERATED IDENTITY AUTHENTICATION AND DUAL-INTERFACE MANAGEMENT ARCHITECTURE}
\end{center}
\vspace{0.5in}
\section*{APPROVED BY}
\vspace{1in}
\noindent\rule{10cm}{0.4pt} \\
Nguyen Van Sinh, Dr.
\vspace{1in}
\noindent\rule{10cm}{0.4pt} \\
Ly Tu Nga, Dr.

\newpage
% ------------------------------------------------------------------ %
% ACKNOWLEDGMENTS -- placeholder; personal, please edit before submitting
% ------------------------------------------------------------------ %
\section*{Acknowledgments}
\addcontentsline{toc}{section}{Acknowledgments}

I would like to express my sincere gratitude to my supervisor, Dr.\
Nguyen Van Sinh, for his guidance, patience, and constructive feedback
throughout the design and development of this thesis. His insights on
system architecture and security practice were instrumental in shaping
the final direction of this project.

I am also grateful to the faculty of the School of Computer Science and
Engineering at the International University for the academic foundation
that made this work possible, and to my family and friends for their
continued support throughout my studies.

\emph{(This section is a placeholder -- please personalize it with the
specific people, mentors, or organizations you want to thank before
submitting.)}

% ------------------------------------------------------------------ %
% TABLE OF CONTENTS / LIST OF TABLES / LIST OF FIGURES
% ------------------------------------------------------------------ %
\newpage
\tableofcontents
\listoftables
\listoffigures

% ------------------------------------------------------------------ %
% ABSTRACT
% ------------------------------------------------------------------ %
\newpage
\section*{Abstract}
\addcontentsline{toc}{section}{Abstract}

Universities increasingly rely on Microsoft 365 as the backbone of
institutional identity and collaboration, yet most learning management
systems still treat authentication as an afterthought bolted onto a
separate credential store. This thesis presents \textbf{myIU}, an
integrated university student portal that delegates authentication
entirely to Microsoft Azure Active Directory through OAuth 2.0 /
OpenID Connect, while keeping full ownership of academic data --
courses, timetables, attendance, assignments, grades, and a lightweight
academic social feed -- inside a purpose-built relational schema.

The system is organized as two independent Spring Boot 3.4 / Java 21
backends sharing a single PostgreSQL 16 database: a \emph{Portal}
backend that owns the schema through Flyway migrations and serves
students and lecturers, and a separate \emph{Admin Panel} backend that
connects to the same database with schema mutation disabled, serving
system administrators. Both issue their own stateless JWTs; sessions can
be revoked instantly through an in-memory blacklist backed by a
persistent \texttt{login\_sessions} table, which also opportunistically
resolves an approximate login location via IP geolocation and, with the
user's explicit consent, a more precise browser GPS lookup reverse
-geocoded through OpenStreetMap Nominatim.

Beyond the core academic workflows, the thesis documents a series of
concrete engineering decisions and the real defects uncovered while
building the system: an IP-spoofing vulnerability in the rate-limiting
and audit-logging path, a foreign key silently missing since the
project's early schema, and an asynchronous transaction bug that left
every login session's location permanently unresolved for weeks without
producing a single visible error. Rather than presenting only the
finished feature set, Chapter 5 evaluates the system against its own
test suite and linters as they stand today, reporting exactly what
passes, what does not, and what has simply never been measured --
in the same spirit of verified, evidence-based reporting that this
thesis's structure was modeled on.

The frontend is built with React 18, Vite, and React Query v5 for
Portal and Admin clients alike, communicating over REST with real-time
push delivered through STOMP over WebSocket. The result is a working,
security-conscious academic platform that demonstrates how federated
identity, a dual-backend administrative boundary, and disciplined
migration-as-code can be combined for a mid-sized university deployment
without resorting to a heavyweight commercial LMS.

\chapter{Introduction}

\section{Background}

Digital transformation in higher education has accelerated the demand
for a single, coherent platform that supports students, lecturers, and
administrators across the full academic lifecycle. Established learning
management systems such as Moodle, Canvas, and Blackboard provide rich
course-management functionality, yet they were largely designed before
institutional identity was centered on a cloud identity provider such
as Microsoft Entra ID (Azure Active Directory). Retrofitting federated
single sign-on onto these platforms is possible but rarely first-class,
and the resulting experience still frequently asks the end user to
maintain a second, LMS-specific credential.

This thesis originates from the practical reality of an international
university in Vietnam where students and lecturers are already issued
Microsoft accounts for the entirety of their academic and administrative
activity. Requiring users to additionally register and remember a
separate account for a learning platform introduces avoidable friction,
encourages password reuse across systems, and complicates account
provisioning whenever a student enrolls or withdraws. The myIU system is
designed to remove this friction by delegating identity verification
entirely to Microsoft Azure AD through the OAuth 2.0 / OpenID Connect
protocol, while every piece of academic data -- courses, enrollment,
schedules, assignments, and grades -- remains fully owned and controlled
by the institution's own database.

\section{Problem Statement}

Academic administration at many Vietnamese universities remains
fragmented across disconnected tools: spreadsheets for grading, email
threads for announcements, and, at best, a legacy web portal for
viewing timetables. Where a unified student information system does
exist, it is typically a monolithic application that conflates
student-facing and administrator-facing concerns in a single codebase
and a single set of credentials, which enlarges the attack surface
exposed to the public internet and makes independent scaling or
hardening of the administrative surface difficult.

There is, in short, a lack of a platform that simultaneously offers:
(1) modern federated authentication against an institution's existing
identity provider, (2) integrated academic workflows -- course
management, timetabling, attendance, assignment submission, and
grading -- inside one system rather than several disconnected tools,
and (3) a clear architectural separation between the interface used by
students and lecturers and the interface used by system administrators,
so that the administrative surface can be deployed, secured, and scaled
independently.

\section{Scope and Objectives}

This thesis focuses on the design and implementation of \textbf{myIU},
a two-application academic platform consisting of a student/lecturer
Portal and a separate Admin Panel. The concrete objectives pursued are:

\begin{enumerate}
\item Design and implement federated authentication integrating
Microsoft Azure AD via OAuth 2.0 / OpenID Connect, allowing sign-in with
institution-issued Microsoft accounts and no separately stored password
for SSO users.
\item Design and build a relational database that serves the core
academic operations of the platform: course and course-group
management, timetabling, attendance, assignment submission, and
grading.
\item Develop a RESTful API using Spring Boot 3.4 and Java 21, applying
a stateless security architecture (JWT) together with modern hardening
practices (rate limiting, CORS policy, revocable sessions).
\item Build a responsive user interface with React 18, Vite, and React
Query v5, cleanly separating server state from local UI state.
\item Design and implement a separate administrative system (Admin
Panel) with its own backend, sharing the Portal's database under a
strict schema-ownership boundary, so that administrative operations
never require exposing the Portal's authentication surface.
\item Implement session-level login tracking, including an opt-in
precise-location feature that reverse-geocodes browser GPS coordinates
when the end user explicitly grants permission, layered on top of a
best-effort IP-based fallback.
\end{enumerate}

\section{Assumptions and Solution}

The design assumes that the deploying institution already operates
Microsoft 365 / Azure AD as its identity provider for staff and
students -- a reasonable assumption for the target context and one that
allows the system to avoid implementing and storing its own password
database for the majority of user accounts.

The proposed solution consists of two independent Spring Boot
backends sharing one PostgreSQL 16 database: the \textbf{Portal}
backend (port 8080), which owns the database schema through Flyway
migrations and serves students and lecturers via Microsoft SSO; and the
\textbf{Admin Panel} backend (port 8081), which connects to the same
database with schema mutation disabled (\texttt{ddl-auto: none}) and
serves system administrators through an entirely separate,
email/password-based authentication flow. Both issue their own
stateless JSON Web Tokens with independent signing secrets, so that a
compromise of one application's token infrastructure cannot be used to
forge access to the other.

\section{Structure of the Thesis}

The remainder of this thesis is organized into five further chapters.
\textbf{Chapter 2: Literature Review} examines existing learning
management systems and enterprise collaboration platforms, together
with the theoretical foundations of federated identity and stateless
API design, to identify the gap that myIU addresses. \textbf{Chapter 3:
Methodology} describes the system's architecture, database design,
class structure, use cases, authentication and authorization mechanism,
session and geolocation handling, real-time notification system, and
migration and configuration strategy. \textbf{Chapter 4: Implementation
and Results} details how the backend, frontend, and admin panel were
actually built, with concrete code excerpts. \textbf{Chapter 5:
Discussion and Evaluation} evaluates the system honestly against its own
test suite and static analysis tooling as they exist today, and walks
through several real defects that were found and fixed during
development. Finally, \textbf{Chapter 6: Conclusion and Future Work}
summarizes the contributions of this thesis and outlines the priorities
for future development.

\chapter{Literature Review}

This chapter reviews the theoretical foundations and existing systems
relevant to the design of myIU. It begins by establishing the
architectural and security concepts that underpin the platform --
single-page application design, federated identity via OAuth 2.0 and
OpenID Connect, stateless JWT-based sessions, role-based access control,
RESTful API design, and migration-as-code -- before examining four
categories of existing solutions: general-purpose learning management
systems, enterprise collaboration suites with academic add-ons, and
locally deployed university portals in Vietnam. The chapter concludes
by identifying the gaps that motivate the design decisions described in
Chapter 3.

\section{Theoretical Foundations}

A single-page application (SPA) architecture, in which the client
renders views dynamically and communicates with the backend exclusively
through a JSON API, was chosen for both the Portal and Admin
frontends. This separates presentation concerns entirely from business
logic, allows the backend to remain a pure API server, and permits the
same backend to eventually serve additional clients (such as a mobile
application) without modification.

Federated identity is realized through OAuth 2.0 (RFC 6749) [1], an
authorization framework that allows a third-party application to obtain
limited access to a user's resources without ever handling that user's
password, and OpenID Connect (OIDC) [2], an identity layer built on top
of OAuth 2.0 that adds a standardized \emph{ID Token} -- itself a JSON
Web Token (JWT, RFC 7519) [3] -- carrying verifiable claims about the
authenticated user. In the Microsoft Azure AD context, the
Authorization Code flow is used: after the user authenticates directly
with Microsoft, the application exchanges an authorization code for an
access token and an ID token, then uses the claims inside the ID token
(email, display name, tenant identifier) to look up or provision a
corresponding local account, entirely decoupling Microsoft's
authentication scope from the application's own authorization scope.

Once a local identity is established, myIU issues its own JWT for every
subsequent request, rather than relying on server-side session state.
This stateless design avoids the operational cost of a shared session
store and scales horizontally without sticky sessions, at the cost of
needing an explicit mechanism to support token revocation (addressed in
Section 3.6). Authorization within the API is enforced through
role-based access control (RBAC), where each authenticated principal
carries one or more roles (\texttt{student}, \texttt{lecturer},
\texttt{admin}) that gate access to specific endpoints and operations.

Finally, the database schema itself is versioned as code using Flyway
[10], a migration tool that applies numbered, immutable SQL scripts in
order and records which have already run. This "migration-as-code"
discipline is what allows two independently deployed backends (Portal
and Admin) to safely agree on a single evolving schema, as discussed
further in Section 3.9.

\section{Moodle}

Moodle is the most widely deployed open-source learning management
system, reporting well over 300 million registered users across its
global installations [6]. Its principal strength is an extensive plugin
ecosystem and deep customizability, which has made it the default choice
for many public universities.

\subsection{Accomplishments}

Moodle's open architecture allows institutions to self-host, audit, and
extend the platform freely, with no licensing cost for the core
product. Its plugin marketplace covers everything from plagiarism
detection to video conferencing, and its course structure (topics,
activities, resources) is flexible enough to model almost any pedagogical
approach.

\subsection{Limitations}

Federated single sign-on with an enterprise identity provider such as
Azure AD is not native to Moodle; it requires installing and
maintaining a third-party SSO plugin, which introduces an additional
dependency and upgrade risk. Moodle's user interface, largely unchanged
in its fundamentals for over a decade, is frequently criticized as
dated and unintuitive for students used to modern consumer applications.
Furthermore, Moodle has no first-class concept of a separate
administrative application boundary: site administrators, course
managers, and end users all operate within the same monolithic
application surface.

\section{Canvas LMS}

Canvas, developed by Instructure, is a commercially hosted LMS that
offers a substantially more modern user interface than Moodle and
supports SSO integration with Microsoft 365 out of the box for
Enterprise-tier customers [7].

\subsection{Accomplishments}

Canvas provides a clean, mobile-friendly interface, a documented REST
API, and native integration paths for SSO and gradebook synchronization
with a Student Information System (SIS), making it attractive for
institutions that can absorb its licensing cost.

\subsection{Limitations}

Canvas is a closed-source, cloud-hosted product; institutions cannot
self-host or deeply customize its internals, and pricing scales with
enrollment, which can be prohibitive for small-to-medium institutions.
Like Moodle, Canvas is architected as a single application rather than
as separate student-facing and administrator-facing systems with
independent security boundaries.

\section{Blackboard Learn}

Blackboard Learn is an enterprise LMS widely deployed at large
universities, offering an extensive feature set spanning course
delivery, analytics, and accessibility tooling.

\subsection{Accomplishments}

Blackboard's maturity and enterprise support contracts make it a
low-risk choice for large institutions with dedicated IT staff, and its
analytics and accessibility compliance tooling are considerably more
developed than most open-source alternatives.

\subsection{Limitations}

Blackboard's feature breadth comes with substantial licensing and
operational cost, and its configuration complexity is frequently
reported as excessive for small-to-medium universities that only need a
subset of its capabilities -- effectively an over-engineered solution
relative to the actual requirements of many deployments.

\section{Microsoft Teams for Education}

Microsoft Teams for Education integrates natively with Azure AD, since
it is itself a Microsoft 365 product, and provides class-team
collaboration features (assignments, channels, meetings) directly
inside the Microsoft ecosystem.

\subsection{Accomplishments}

Because Teams shares the same identity plane as the rest of an
institution's Microsoft 365 tenant, single sign-on is entirely native
with zero additional configuration, and the collaboration and video
conferencing tooling is best-in-class.

\subsection{Limitations}

Teams for Education is positioned as a collaboration tool rather than a
purpose-built academic information system: it has no native concept of
a structured timetable, no daily attendance tracking, and no
component-based grade breakdown. Institutions that adopt it as their
primary academic platform must still build or bolt on separate tooling
for these core academic record-keeping needs.

\section{Existing Vietnamese University Student Portals}

Several public Vietnamese universities, including large technical
universities in Ho Chi Minh City, operate in-house student information
portals for viewing timetables, grades, and administrative
announcements.

\subsection{Accomplishments}

These portals are tailored to local academic regulations and
administrative processes (credit-based grading scales, local
semester structures, Vietnamese-language interfaces) in ways that
international commercial products rarely are out of the box.

\subsection{Limitations}

Many of these portals were built years ago on legacy web stacks, have
dated interfaces, lack a companion mobile application, and predate
modern federated identity standards -- authentication is typically a
locally stored username/password pair rather than an integration with
the institution's existing identity provider, reintroducing exactly the
credential-duplication problem described in Section 1.1.

\section{Gap Identification and Proposed Enhancements}

Table \ref{tab:lit-comparison} summarizes the comparison across the
five reviewed categories of existing solutions against myIU on the
dimensions most relevant to this thesis's objectives.

\begin{longtable}{@{}p{0.20\linewidth}p{0.15\linewidth}p{0.15\linewidth}p{0.13\linewidth}p{0.15\linewidth}@{}}
\caption{Comparison of Existing Platforms and myIU}
\label{tab:lit-comparison}\\
\toprule
\textbf{Criterion} & \textbf{Moodle} & \textbf{Canvas / Blackboard} & \textbf{MS Teams Edu} & \textbf{myIU} \\
\midrule
\endhead
Microsoft Azure AD SSO & Third-party plugin & Native (Enterprise tier) & Native & Native (OIDC) \\
Timetable management & Add-on module & Limited & None & Full CRUD \\
Account provisioning & CSV / plugin & API / SIS sync & AD-synced & Manual, per-account, via Admin UI \\
Customizability & High (self-hosted PHP) & Low--Medium & Low & High (open codebase) \\
API architecture & REST / web services & REST & Microsoft Graph API & REST + stateless JWT \\
Separate admin surface & No & No & No & Yes (independent backend) \\
\bottomrule
\end{longtable}

From this comparison, three gaps motivate the design of myIU: (1) the
absence of a highly customizable platform that natively integrates
enterprise Microsoft SSO in a way tailored to the Vietnamese university
context; (2) the lack of an open solution that combines timetabling,
attendance, assignment submission, and an academic social feed inside a
single coherent system, rather than several disconnected tools; and (3)
the relative rarity, in the Spring Boot ecosystem specifically, of a
documented dual-backend architecture that cleanly separates the
end-user-facing interface from the administrative interface behind an
independent security boundary.

\section{Summary}

The systems reviewed in this chapter demonstrate that mature learning
management platforms exist and are effective at their core mission, but
each carries trade-offs that make them a poor fit for a
mid-sized, Microsoft-365-centric institution seeking a highly
customizable, security-conscious, and administratively separated
platform: Moodle and Canvas are either difficult to integrate natively
with enterprise SSO or expensive to license; Blackboard is
over-engineered for the target scale; Microsoft Teams for Education
lacks core academic record-keeping features; and existing local
Vietnamese portals predate modern identity federation entirely. myIU is
designed specifically to close these gaps, as detailed in the following
chapter.

\chapter{Methodology}

\section{Overview}

myIU is built as two independent Spring Boot 3.4.1 applications on
Java 21 -- enabling Project Loom virtual threads and modern language
features such as records -- sharing a single PostgreSQL 16 database.
The \textbf{Portal} backend serves students and lecturers and is the
sole owner of the database schema, evolved exclusively through Flyway
migrations. The \textbf{Admin Panel} backend serves system
administrators, connects to the same database, but is configured with
schema mutation disabled so that it can never drift the schema that the
Portal owns. Both frontends are built with React 18 and Vite; the
Portal frontend targets students and lecturers, while the Admin
frontend is a completely separate single-page application targeting
administrators only. This section introduces the architectural
rationale behind that split before the remaining sections detail the
database, class structure, use cases, security mechanisms, and
operational configuration in turn. Figure \ref{fig:dfd0} gives a
context-level view of the whole system as a single black box together
with the external actors and services it exchanges data with, before
the following sections decompose it further.

% SCREENSHOT SOURCE: docs/system-design.html, diagram 4.3
% "DFD Level 0 -- Context Diagram" (article id="s4-3")
\begin{figure}[htbp]
\centering
\includegraphics[width=0.9\linewidth]{diagram-4-3-dfd-level0.png}
\caption{DFD Level 0 -- Context Diagram (source: \texttt{docs/system-design.html}, diagram 4.3)}
\label{fig:dfd0}
\end{figure}

\section{Architecture Overview and Design Rationale}

\subsection{Design Rationale}

The technology choices for myIU were driven by three requirements:
strict separation between end-user and administrative concerns,
predictable behavior under concurrent academic-term traffic spikes
(course registration, grade release), and a schema that can evolve
safely as new academic features are added.

\textbf{Layered architecture.} Both backends follow a conventional
Controller~$\rightarrow$~Service~$\rightarrow$~Repository~$\rightarrow$~Entity
layering. Controllers are kept thin and translate HTTP requests into
service calls; services hold business logic and transaction boundaries;
Spring Data JPA repositories provide typed data access; and JPA
entities map directly onto the underlying PostgreSQL tables.

\textbf{Dual-system architecture.} The Portal backend is the exclusive
owner of the database schema: only its Flyway configuration is allowed
to apply migrations (V1 through V12 at the time of writing). The Admin
backend is configured with \texttt{spring.jpa.hibernate.ddl-auto: none},
so Hibernate never attempts to alter the schema on that side -- a
convention-based boundary rather than a database-permission-enforced
one, discussed further as a limitation in Section 5.2.

\textbf{Iterative refinement as an architectural decision.} An early
version of the system implemented bulk user provisioning through an
uploaded Excel spreadsheet, processed with Apache POI, in both
backends. A full audit of the Portal codebase later revealed that its
copy of this feature was entirely dead code: a
\texttt{UserProvisioningService} existed, but no controller ever
injected or called it, meaning the feature had never actually been
reachable through the Portal API. Rather than deleting only the
unreachable code, the team removed the feature end to end from both
backends: the only endpoint that was genuinely reachable
(\texttt{POST /api/users/provision} on the Admin side) increased attack
surface for a rarely used operation, the Apache POI dependency existed
solely to support it, and creating or deactivating accounts one at a
time through the Admin UI was judged sufficient for the system's actual
scale. This is presented as a deliberate architectural decision, not
merely a bug fix: choosing to remove a feature is itself a design
choice with a rationale, and the audit trail columns
(\texttt{provisioned\_by}, \texttt{provisioned\_at}) were deliberately
kept on the \texttt{users} table because they retain value independent
of how an account was actually created.

\subsection{Layers Architecture}

Figure \ref{fig:architecture} shows the container-level architecture of
the whole system. The two applications are fully independent processes
communicating with clients over REST and sharing only the PostgreSQL
instance; there is no shared code or shared runtime between Portal and
Admin.

% SCREENSHOT SOURCE: docs/system-design.html, diagram 2.1
% "Architecture Diagram -- C4 Container Level" (article id="s2-1")
\begin{figure}[htbp]
\centering
\includegraphics[width=0.9\linewidth]{diagram-2-1-architecture-c4.png}
\caption{Architecture Diagram -- C4 Container Level (source: \texttt{docs/system-design.html}, diagram 2.1)}
\label{fig:architecture}
\end{figure}

Figure \ref{fig:dfd1} complements this container view with a
process-level decomposition of the Portal backend into seven
functional service clusters and their shared data stores.

% SCREENSHOT SOURCE: docs/system-design.html, diagram 4.4
% "DFD Level 1 -- Process Decomposition" (article id="s4-4")
\begin{figure}[htbp]
\centering
\includegraphics[width=0.9\linewidth]{diagram-4-4-dfd-level1.png}
\caption{DFD Level 1 -- Process Decomposition (source: \texttt{docs/system-design.html}, diagram 4.4)}
\label{fig:dfd1}
\end{figure}

Advantages of this layering, common to both backends individually,
include separation of concerns (each layer has a single
responsibility, easing testing and maintenance), dependency direction
(services depend on repository interfaces rather than concrete data
access code), and testability (each layer can be exercised in
isolation with mocks). At the system level, the dual-backend split
additionally minimizes the attack surface reachable from the public
internet -- the Admin backend can be deployed on an internal network
segment entirely inaccessible to end users -- and allows the Portal,
which receives far more traffic, to be scaled independently of the
lower-traffic Admin backend.

\section{Database Design}

The PostgreSQL schema is organized into five domain clusters, all owned
by the Portal's Flyway migrations. Every table uses a \texttt{UUID}
primary key generated by PostgreSQL's \texttt{gen\_random\_uuid()}
function rather than an application-side generator, keeping key
generation consistent regardless of which backend performs the insert.

A cross-cutting design decision concerns referential integrity on
foreign keys that reference \texttt{users(id)}. Seven such foreign keys
exist across the schema, and they fall into two distinct categories by
nature: (a) \emph{data owned exclusively by that user}
(\texttt{form\_submissions.submitter\_id},
\texttt{support\_tickets.submitter\_id}), where cascading a delete is
unambiguously safe, and (b) \emph{content a user created that other
people depend on} (\texttt{courses.creator\_id},
\texttt{course\_groups.lecturer\_id}, \texttt{course\_posts.author\_id},
\texttt{course\_grades.graded\_by}, \texttt{form\_templates.created\_by}),
where cascading means deleting one user can also delete a course and
every other student's grades inside it. The team's final decision --
cascading both categories -- was made deliberately for the system's
current demo/seed-data context, with an explicit note in the technical
documentation that this choice should be revisited before the system is
operated on real student data (see Section 5.2).

\subsection{Identity and Access Domain}

This cluster covers user accounts, administrative accounts, and login
session tracking, including the opt-in precise-location columns added
in migration V12.

\begin{longtable}{@{}p{0.20\linewidth}p{0.34\linewidth}p{0.36\linewidth}@{}}
\caption{Key Entities in the Identity and Access Domain}\\
\toprule
\textbf{Entity} & \textbf{Key Attributes} & \textbf{Design Rationale} \\
\midrule
\endhead
users & email (unique), username (unique), password (nullable), auth\_provider, ms\_tenant\_id, is\_active, provisioned\_by, provisioned\_at &
  \texttt{password} is null for Microsoft SSO accounts; \texttt{is\_active} lets an admin deactivate an account without deleting its history \\
user\_roles & user\_id (FK), role &
  Element-collection table backing \texttt{Set<String> roles}; eagerly fetched since Spring Security needs roles on every authenticated request \\
admin\_users & email (unique), password (BCrypt) &
  Deliberately separate from \texttt{users}; shared only at the database level with myIU-admin, never through the Portal's authentication flow \\
login\_sessions & ip\_address, country, city, latitude, longitude, province, district, ward, browser, os, device\_type, is\_revoked &
  Latitude/longitude/province/district/ward (migration V12) are nullable and populated only if the user opts in to browser GPS; otherwise the row is populated from best-effort IP geolocation only \\
\bottomrule
\end{longtable}

\subsection{Academic Core Domain}

This cluster models the structural relationships between courses,
their sections, and enrolled students.

\begin{longtable}{@{}p{0.20\linewidth}p{0.34\linewidth}p{0.36\linewidth}@{}}
\caption{Key Entities in the Academic Core Domain}\\
\toprule
\textbf{Entity} & \textbf{Key Attributes} & \textbf{Design Rationale} \\
\midrule
\endhead
courses & name, code, semester, cover\_color, is\_active, creator\_id (FK) &
  \texttt{cover\_color} drives the visual identity of the course card and timetable block in the frontend \\
course\_groups & course\_id (FK), lecturer\_id (FK), name, max\_students, room &
  A course can be sectioned into multiple groups, each led by a (possibly different) lecturer \\
group\_members & group\_id (FK), course\_id (FK), student\_id (FK), UNIQUE(group\_id, student\_id) &
  \texttt{course\_id} is deliberately denormalized alongside \texttt{group\_id} so that "which courses is this student in" is a single-table lookup, without joining through \texttt{course\_groups} \\
course\_schedules & course\_id (FK), day\_of\_week, start\_time, end\_time, room &
  Backs the weekly timetable grid; a CHECK constraint enforces \texttt{end\_time > start\_time} \\
\bottomrule
\end{longtable}

\subsection{Assessment and Progress Domain}

This cluster tracks the graded, day-to-day academic record of each
student within a course.

\begin{longtable}{@{}p{0.20\linewidth}p{0.34\linewidth}p{0.36\linewidth}@{}}
\caption{Key Entities in the Assessment and Progress Domain}\\
\toprule
\textbf{Entity} & \textbf{Key Attributes} & \textbf{Design Rationale} \\
\midrule
\endhead
course\_posts & course\_id (FK), group\_id (FK), author\_id (FK), title, content, type, due\_date &
  A single model backs both plain announcements and gradable assignment posts, distinguished by \texttt{type} \\
assignment\_submissions & course\_post\_id, student\_id (FK), file\_url, score, graded\_by (FK), status &
  \texttt{course\_post\_id} originally had no foreign key at all -- a real defect discovered and fixed during this project, discussed in Section 5.2 \\
attendance\_records & course\_id (FK), student\_id (FK), date, status, marked\_by (FK), UNIQUE(course\_id, student\_id, date) &
  The unique constraint prevents a student from being marked twice for the same session \\
course\_grades & course\_id (FK), student\_id (FK), quiz, exercise, lab, midterm, project, final\_score &
  Component scores are stored individually rather than only a computed total, preserving the full grade breakdown for review \\
\bottomrule
\end{longtable}

Figure \ref{fig:erd-academic} shows the full entity-relationship
diagram covering the Identity, Academic Core, and Assessment and
Progress domains introduced above.

% SCREENSHOT SOURCE: docs/system-design.html, diagram 4.1
% "ERD -- Core Academic Domain" (article id="s4-1")
\begin{figure}[htbp]
\centering
\includegraphics[width=0.95\linewidth]{diagram-4-1-erd-academic.png}
\caption{Entity-Relationship Diagram -- Core Academic Domain (source: \texttt{docs/system-design.html}, diagram 4.1)}
\label{fig:erd-academic}
\end{figure}

\subsection{Communication and Social Domain}

This cluster supports the academic social feed and system
notifications.

\begin{longtable}{@{}p{0.20\linewidth}p{0.34\linewidth}p{0.36\linewidth}@{}}
\caption{Key Entities in the Communication and Social Domain}\\
\toprule
\textbf{Entity} & \textbf{Key Attributes} & \textbf{Design Rationale} \\
\midrule
\endhead
posts & author\_id (FK), caption, location, likes\_count &
  A lightweight, Instagram-style feed for academic and campus announcements shared across the student body \\
post\_media / comments / saves & post\_id (FK), url/type or body, user\_id (FK) &
  Standard supporting tables for a social feed: attached media, threaded comments, and per-user bookmarking \\
notifications & user\_id (FK), type, title, message, is\_read, related\_id &
  \texttt{type} (\texttt{grade}, \texttt{form\_approved}, \texttt{course}, \texttt{system}, ...) drives both the icon shown in the UI and the WebSocket topic used for real-time delivery (Section 3.8) \\
\bottomrule
\end{longtable}

\subsection{Administrative Services Domain}

This cluster covers dynamic form-based administrative requests and
technical support.

\begin{longtable}{@{}p{0.20\linewidth}p{0.34\linewidth}p{0.36\linewidth}@{}}
\caption{Key Entities in the Administrative Services Domain}\\
\toprule
\textbf{Entity} & \textbf{Key Attributes} & \textbf{Design Rationale} \\
\midrule
\endhead
form\_templates & title, fields (JSON), is\_active, created\_by (FK) &
  \texttt{fields} stores a JSON schema so new administrative form types (scholarship application, exemption request, ...) can be added without a code change or migration \\
form\_submissions & template\_id (FK), submitter\_id (FK), data (JSON), status, reviewed\_by (FK) &
  \texttt{status} transitions \texttt{pending} $\rightarrow$ \texttt{approved}/\texttt{rejected}; a status change asynchronously sends an email to the submitter if SMTP is configured \\
support\_tickets & submitter\_id (FK), service, need, description, status, response, resolved\_by (FK) &
  Simple ticketing model for general technical support requests, separate from the structured form-submission workflow \\
stored\_files & original\_name, stored\_name, content\_type, size, uploaded\_by (FK) &
  Backs a single shared upload endpoint used by assignment submissions, form attachments, and post media alike \\
\bottomrule
\end{longtable}

Figure \ref{fig:erd-platform} shows the entity-relationship diagram for
the Communication and Social, and Administrative Services domains,
together with the \texttt{login\_sessions} table introduced in Section
3.3.1.

% SCREENSHOT SOURCE: docs/system-design.html, diagram 4.2
% "ERD -- Platform & Support Domain" (article id="s4-2")
\begin{figure}[htbp]
\centering
\includegraphics[width=0.95\linewidth]{diagram-4-2-erd-platform.png}
\caption{Entity-Relationship Diagram -- Platform and Support Domain (source: \texttt{docs/system-design.html}, diagram 4.2)}
\label{fig:erd-platform}
\end{figure}

\section{Class Diagram Design}

The domain model is implemented with Jakarta Persistence API (JPA)
entities and Lombok annotations to minimize boilerplate. Figure
\ref{fig:classes} summarizes the core academic entities and their
relationships, mirroring the tables introduced in Section 3.3.

% SCREENSHOT SOURCE: docs/system-design.html, diagram 3.1
% "Class Diagram -- Core Academic Domain" (article id="s3-1")
\begin{figure}[htbp]
\centering
\includegraphics[width=0.95\linewidth]{diagram-3-1-class-diagram.png}
\caption{Class Diagram -- Core Academic Domain (source: \texttt{docs/system-design.html}, diagram 3.1)}
\label{fig:classes}
\end{figure}

Beyond the entity layer, each domain area is organized into a
Repository (Spring Data JPA), a Service holding transactional business
logic, a set of DTOs used for API projection so that internal entity
structure is never leaked directly to the client, and a thin
Controller. Table \ref{tab:services} summarizes the key services and
their responsibilities.

\begin{longtable}{@{}p{0.24\linewidth}p{0.66\linewidth}@{}}
\caption{Key Backend Services and Responsibilities}
\label{tab:services}\\
\toprule
\textbf{Service} & \textbf{Responsibilities} \\
\midrule
\endhead
JwtService & JWT issuance and validation, blacklist lookup for revoked tokens \\
SessionService & Login session creation, heartbeat, revocation, IP resolution \\
GeoIpService & Asynchronous best-effort IP geolocation and opt-in GPS reverse-geocoding \\
TimetableService & Aggregates course schedules into a per-user weekly timetable \\
NotificationService & Persists notifications and pushes them over STOMP/WebSocket \\
StorageService & Stores uploaded files and issues retrievable URLs \\
UserProvisioningService & Legacy Excel-based bulk provisioning (removed end-to-end; see Section 3.2.1) \\
\bottomrule
\end{longtable}

\section{Use Case Analysis}

The system serves three internal actors -- Student, Lecturer, and
Administrator -- plus one external actor, Microsoft Azure AD, which
participates in the authentication use case as an identity provider.

\subsection{Actor Descriptions}

\begin{enumerate}
\item \textbf{Student}: Signs in via Microsoft SSO, views their
personal timetable, browses enrolled courses and course materials,
submits assignments, checks attendance and grades, participates in the
academic social feed, submits administrative forms and support
tickets, and manages their own login sessions.
\item \textbf{Lecturer}: Signs in via the same Microsoft SSO flow as
students but is bound to a \texttt{lecturer} role; creates and manages
course groups, publishes course posts (including gradable
assignments), grades submissions, marks attendance, and participates in
the social feed.
\item \textbf{Administrator}: Authenticates with a locally stored
email/password pair through the entirely separate Admin Panel; manages
user accounts, courses, and grades; moderates content; handles support
tickets and form-submission reviews; broadcasts notifications; and
monitors and revokes login sessions, including their resolved
geolocation.
\end{enumerate}

Figure \ref{fig:usecase-diagram} shows all three actors together with
the full set of use cases they participate in, spanning both the
Portal and the Admin Panel systems.

% SCREENSHOT SOURCE: docs/system-design.html, diagram 1.4
% "Use Case Diagram -- Full System" (article id="s1-4")
\begin{figure}[htbp]
\centering
\includegraphics[width=0.95\linewidth]{diagram-1-4-use-case.png}
\caption{Use Case Diagram -- Full System (source: \texttt{docs/system-design.html}, diagram 1.4)}
\label{fig:usecase-diagram}
\end{figure}

\subsection{Use Case Groups}

\begin{itemize}
\item \textbf{Authentication}: Microsoft SSO sign-in (Student/Lecturer),
Admin email/password sign-in, JWT logout, session revocation.
\item \textbf{Course Management}: create/update/delete course, manage
course groups and enrollment.
\item \textbf{Timetable}: view personal timetable, add/edit/delete
schedule entries.
\item \textbf{Assignment Submission and Grading}: publish assignment
post, submit assignment, grade submission.
\item \textbf{Attendance}: mark single or bulk attendance, view
attendance history.
\item \textbf{Grade Lookup}: view course grades, view own grades across
all courses.
\item \textbf{Social Feed}: create/view/like/comment/save posts.
\item \textbf{Form Submission}: submit a dynamic administrative form,
review and approve/reject a submission.
\item \textbf{Admin Dashboard}: platform-wide statistics, user
management, content moderation, notification broadcast.
\item \textbf{Session Monitoring with Location}: view active sessions,
revoke a session, opt in to precise GPS-based location resolution.
\end{itemize}

Table \ref{tab:usecases} lists a representative subset of use cases
together with their corresponding API endpoint, drawn directly from the
system's route reference.

\begin{longtable}{@{}p{0.09\linewidth}p{0.24\linewidth}p{0.38\linewidth}p{0.17\linewidth}@{}}
\caption{Key Use Cases with API Endpoints}
\label{tab:usecases}\\
\toprule
\textbf{ID} & \textbf{Use Case} & \textbf{API Endpoint} & \textbf{Actor} \\
\midrule
\endhead
UC-01 & Microsoft SSO Sign-in & GET /oauth2/authorization/microsoft & Student/Lecturer \\
UC-02 & Get Current User & GET /api/auth/me & Student/Lecturer \\
UC-03 & List/View Courses & GET /api/courses, GET /api/courses/:id & All \\
UC-04 & Create Course & POST /api/courses & Lecturer \\
UC-05 & Manage Course Groups & GET/POST /api/courses/:id/groups & Lecturer \\
UC-06 & Manage Group Members & GET/POST/DELETE /api/groups/:id/members & Lecturer \\
UC-07 & View Timetable & GET /api/timetable & Student/Lecturer \\
UC-08 & Manage Timetable Entry & POST/PUT/DELETE /api/timetable & Lecturer/Admin \\
UC-09 & Submit Assignment & POST /api/submissions & Student \\
UC-10 & Grade Assignment & PUT /api/submissions/:id/grade & Lecturer \\
UC-11 & Mark Attendance & POST /api/attendance, POST /api/attendance/bulk & Lecturer \\
UC-12 & View Own Attendance & GET /api/attendance/mine & Student \\
UC-13 & View Course Grades & GET /api/courses/:id/grades & All \\
UC-14 & View Own Grades & GET /api/grades/me & Student \\
UC-15 & Social Feed Post & GET/POST /api/posts & All \\
UC-16 & Like/Comment on Post & POST /api/posts/:id/like, /comments & All \\
UC-17 & Submit Administrative Form & POST /api/forms/submit & Student/Lecturer \\
UC-18 & Review Form Submission & PATCH /api/forms/submissions/:id & Admin \\
UC-19 & File Support Ticket & POST /api/support & All \\
UC-20 & Respond to Support Ticket & PUT /api/admin/tickets/:id & Admin \\
UC-21 & Manage User Accounts & GET /api/admin/users, PUT /:id/status & Admin \\
UC-22 & View Platform Stats & GET /api/admin/stats & Admin \\
UC-23 & View/Revoke Login Sessions & GET /api/sessions, DELETE /api/sessions/:id & Student/Lecturer \\
UC-24 & Report Precise Location & PUT /api/sessions/current/location & Student/Lecturer \\
\bottomrule
\end{longtable}

\subsection{Critical Flow: SSO Login to JWT Issuance}

Figure \ref{fig:activity-sso} first presents the business-process view
of Microsoft SSO sign-in, from the moment the user lands on the Portal
to the moment the Home page is rendered.

% SCREENSHOT SOURCE: docs/system-design.html, diagram 1.1
% "Activity Diagram -- Student Authentication (Microsoft SSO)" (article id="s1-1")
\begin{figure}[htbp]
\centering
\includegraphics[width=0.8\linewidth]{diagram-1-1-activity-sso.png}
\caption{Activity Diagram -- Student Authentication via Microsoft SSO (source: \texttt{docs/system-design.html}, diagram 1.1)}
\label{fig:activity-sso}
\end{figure}

Figure \ref{fig:sso-flow} then traces the same flow at the technical
sequence-diagram level, from the moment a user clicks ``Sign in with
Microsoft'' to the moment the frontend holds a valid JWT. After the
browser is redirected to Azure AD and the user authenticates directly
with Microsoft, the backend's OAuth2 callback exchanges the
authorization code for tokens, calls the Microsoft Graph API to
retrieve the user's profile, looks up or creates the corresponding
local \texttt{User} row, creates a \texttt{LoginSession} row, and
finally issues a JWT carrying the user's identifier, roles, and a
unique \texttt{jti} claim used for revocation (Section 3.6).

% SCREENSHOT SOURCE: docs/system-design.html, diagram 3.2
% "Sequence Diagram -- Microsoft SSO Login" (article id="s3-2")
\begin{figure}[htbp]
\centering
\includegraphics[width=0.95\linewidth]{diagram-3-2-sequence-sso.png}
\caption{Sequence Diagram -- Microsoft SSO Login to JWT Issuance (source: \texttt{docs/system-design.html}, diagram 3.2)}
\label{fig:sso-flow}
\end{figure}

\subsection{Critical Flow: Assignment Submission and Grading}

This flow spans the full lifecycle of a gradable assignment: a
lecturer publishes an assignment post, a student submits a file or text
answer, and the lecturer later reviews and grades that submission,
with a real-time notification delivered to each party at the relevant
step (Section 3.8). Figure \ref{fig:activity-assignment} shows the
process view.

% SCREENSHOT SOURCE: docs/system-design.html, diagram 1.2
% "Activity Diagram -- Assignment Submission & Grading Lifecycle" (article id="s1-2")
\begin{figure}[htbp]
\centering
\includegraphics[width=0.8\linewidth]{diagram-1-2-activity-assignment.png}
\caption{Activity Diagram -- Assignment Submission and Grading Lifecycle (source: \texttt{docs/system-design.html}, diagram 1.2)}
\label{fig:activity-assignment}
\end{figure}

Figure \ref{fig:sequence-assignment} shows the corresponding technical
interaction between the student, lecturer, backend, storage service,
and the WebSocket notification channel.

% SCREENSHOT SOURCE: docs/system-design.html, diagram 3.3
% "Sequence Diagram -- Assignment Submission & Grading" (article id="s3-3")
\begin{figure}[htbp]
\centering
\includegraphics[width=0.95\linewidth]{diagram-3-3-sequence-assignment.png}
\caption{Sequence Diagram -- Assignment Submission and Grading (source: \texttt{docs/system-design.html}, diagram 3.3)}
\label{fig:sequence-assignment}
\end{figure}

\subsection{Critical Flow: Form Submission and Admin Approval}

Dynamic administrative forms (scholarship applications, exemption
requests, and similar) follow a submit-then-review workflow ending in
an administrator's approval or rejection decision, which triggers an
outbound notification (and, if configured, an email) back to the
submitter. Figure \ref{fig:activity-forms} shows this workflow in full,
including the branch where a rejected submission records a rejection
reason.

% SCREENSHOT SOURCE: docs/system-design.html, diagram 1.3
% "Activity Diagram -- Form Submission & Admin Approval" (article id="s1-3")
\begin{figure}[htbp]
\centering
\includegraphics[width=0.8\linewidth]{diagram-1-3-activity-forms.png}
\caption{Activity Diagram -- Form Submission and Admin Approval (source: \texttt{docs/system-design.html}, diagram 1.3)}
\label{fig:activity-forms}
\end{figure}

\section{Authentication and Authorization Architecture}

Authentication in the Portal is built on three cooperating components:
\texttt{SecurityConfig}, which wires the Spring Security filter chain
into stateless OAuth2 login; \texttt{JwtService}, which issues and
validates tokens; and \texttt{JwtAuthenticationFilter}, which runs
before Spring Security's default authentication filter and enforces
both signature validity and session-level revocation on every request.
\texttt{SessionCreationPolicy.STATELESS} ensures Spring Security never
creates an \texttt{HttpSession}, consistent with a pure REST API; CORS
is restricted per environment to the known frontend origins.

Each issued JWT carries a unique \texttt{jti} (JWT ID) claim. On
logout, that \texttt{jti} is added to an in-memory blacklist, allowing
immediate revocation on the issuing node without needing to persist
every token ever issued. Algorithm \ref{alg:jwt-validate} summarizes the
validation path executed on every authenticated request.

\begin{algorithm}[htbp]
\caption{JWT Validation with Session Revocation Check}
\label{alg:jwt-validate}
\begin{algorithmic}[1]
\Procedure{ValidateRequest}{$rawToken$}
  \If{$rawToken = null$}
    \State \Return \textsc{Unauthenticated}
  \EndIf
  \State $claims \gets \Call{ParseAndVerifySignature}{rawToken}$
  \If{$claims = null \lor \Call{IsExpired}{claims}$}
    \State \Return \textsc{Error\_Invalid\_Token}
  \EndIf
  \State $jti \gets claims.jti$
  \If{$\Call{JwtBlacklist.IsRevoked}{jti}$}
    \State \Return \textsc{Error\_Token\_Revoked}
  \EndIf
  \State $sessionId \gets claims.sessionId$
  \If{$sessionId \neq null$}
    \State $sessionOk \gets \Call{TryCheckSessionNotRevoked}{sessionId}$
    \Comment{Wraps the repository call in a catch-all; on exception (e.g.\ DB
    temporarily unavailable) it logs a warning and returns \textsc{true} --
    fail-open rather than blocking every request}
    \If{$\lnot sessionOk$}
      \State \Return \textsc{Error\_Session\_Revoked}
    \EndIf
  \EndIf
  \State \Call{SecurityContext.SetAuthentication}{claims}
  \State \Return \textsc{Authenticated}
\EndProcedure
\end{algorithmic}
\end{algorithm}

The session-level check exists because the in-memory blacklist alone
would not survive a restart or be shared across multiple backend
instances; the \texttt{login\_sessions} table's \texttt{is\_revoked}
column gives a persistent, durable second layer of revocation. Notably,
the persistent check is deliberately \emph{fail-open}: if the database
is temporarily unreachable, the request is allowed to proceed rather
than rejected outright, trading a small security margin for
availability during transient database issues.

\subsection{Client IP Trust Boundary}

The \texttt{X-Forwarded-For} and \texttt{X-Real-IP} HTTP headers are
used both as the key for login rate-limiting (to block brute-force
attempts) and as the source IP recorded in \texttt{login\_sessions} for
audit and location purposes. Both headers are entirely controlled by
the client sending the request; trusting them unconditionally would let
any attacker who connects directly to the backend (bypassing any
reverse proxy) present a different fabricated IP on every request,
defeating the rate limiter completely and forging the location shown to
the user. Algorithm \ref{alg:ip-trust} describes the trust boundary
adopted to close this gap.

\begin{algorithm}[htbp]
\caption{Trusted-Proxy IP Resolution}
\label{alg:ip-trust}
\begin{algorithmic}[1]
\Procedure{ExtractClientIp}{$request$}
  \State $peer \gets request.\Call{getRemoteAddr}{}$ \Comment{TCP peer; cannot be spoofed by the client}
  \If{$\Call{IsPrivateOrLoopback}{peer}$}
    \Comment{Request genuinely arrived through an internal reverse proxy}
    \State $forwarded \gets request.\Call{getHeader}{\text{"X-Forwarded-For"}}$
    \If{$forwarded \neq null$}
      \State \Return \Call{FirstAddress}{forwarded}
    \EndIf
  \EndIf
  \State \Return $peer$ \Comment{Untrusted peer: ignore client-supplied headers entirely}
\EndProcedure
\end{algorithmic}
\end{algorithm}

This is the same "trust proxy" pattern used by mainstream production
frameworks (Express's \texttt{trust proxy} setting, Nginx's
\texttt{real\_ip} module), but it was not present in the system's
earlier revisions -- its absence and the concrete impact are discussed
as a discovered vulnerability in Section 5.2.

\section{Federated Identity, Session Management, and Precise Geolocation}

No password is ever stored for a Microsoft SSO account; the
\texttt{password} column on \texttt{users} is simply left null for
those rows. Because the Portal and Admin backends are entirely separate
applications, each signs its JWTs with its own independent secret --
compromising one application's signing key cannot be used to forge a
token accepted by the other.

IP-based geolocation (via ip-api.com) is inherently limited to
city-level precision, since an ISP's allocated IP block is not fine
enough to imply a specific ward or district. To offer a more precise
location in the session-management UI, the system layers an
\emph{opt-in} flow on top: after a successful login, the frontend asks
the browser once for GPS permission; if granted, the coordinates are
sent to the backend and reverse-geocoded through OpenStreetMap's
Nominatim service down to ward/district/province. If permission is
denied, the system simply falls back to the existing IP-based
behavior -- nothing else changes.

Nominatim's address response for Vietnamese locations does not expose a
dedicated "province" field; the province/city appears only as free text
inside \texttt{display\_name}. After verifying against real
coordinates, the team found that, once the postal code and country
name are stripped from the end of the string, the last three
comma-separated segments \emph{usually} correspond, in order, to
[ward/commune, district or provincial city, province/city] -- making
positional parsing more reliable in practice than guessing which OSM
address field corresponds to which Vietnamese administrative level.
However, this "always three levels" assumption is not universally true:
testing with real user coordinates showed that, following recent
administrative mergers, some wards sit directly under the
province/city level with no intervening district tier at all -- for
example, \texttt{display\_name} for one real ward resolves to exactly
two segments ("Phuong Binh Thanh, Ho Chi Minh City"), not three. A naive
fixed-offset parser in that case mistakenly reused the ward's own name
as the "district," producing a visibly duplicated result. Algorithm
\ref{alg:geo-dedup} shows the deduplication step added to fix this: if
the computed district equals either the ward or the province, it is
discarded (treated as absent) rather than displayed as a duplicate --
accepting a slightly less precise result in that edge case in exchange
for never showing obviously wrong, repeated text.

\begin{algorithm}[htbp]
\caption{Reverse-Geocoded Address Deduplication}
\label{alg:geo-dedup}
\begin{algorithmic}[1]
\Procedure{ParseVietnameseAddress}{$displayName$}
  \State $segments \gets \Call{SplitAndTrimTrailingPostalAndCountry}{displayName}$
  \State $ward, district, province \gets \Call{LastThreeSegments}{segments}$
  \If{$district = ward \lor district = province$}
    \State $district \gets null$ \Comment{Discard the duplicate instead of displaying it}
  \EndIf
  \State \Return $(ward, district, province)$
\EndProcedure
\end{algorithmic}
\end{algorithm}

A related, purely internal defect was found in the same code path
during review: the URL-building call used
\texttt{String.format("\%f", lat, lng)}, whose output depends on the
JVM's default locale. On a server configured with a
comma-as-decimal-separator locale, this would silently corrupt the
query string sent to Nominatim (printing \texttt{10,77} instead of
\texttt{10.77}). The bug never manifested on the development machine
(whose JVM defaults to \texttt{en\_US}) but was real, and was fixed
proactively by pinning \texttt{Locale.US} explicitly -- an example of
the value of actively reviewing code rather than waiting for a
locale-dependent bug to surface in production.

\section{Real-Time Notification System}

Rather than polling, notifications are delivered over a real WebSocket
push channel. \texttt{NotificationService.push()} persists a
\texttt{Notification} row and forwards the same payload to a
per-user STOMP topic, \texttt{/topic/notifications/\{userId\}}, which
the frontend subscribes to through \texttt{useNotificationSocket.ts}
using SockJS as a fallback transport for environments where raw
WebSocket connections are blocked. A dedicated
\texttt{WebSocketAuthInterceptor} authenticates the caller's JWT at the
STOMP \texttt{CONNECT} frame itself, before any subscription is
accepted, so that the WebSocket channel carries the same authentication
guarantees as the REST API. Notifications are triggered by: assignment
submission (notifying the lecturer), grading (notifying the student),
form approval or rejection (notifying the submitter, and, if SMTP is
configured, also by email), and administrator broadcasts.

\section{Data Migration Strategy with Flyway}

Flyway is configured with two source locations: \texttt{db/migration},
containing schema-defining migrations that run in every environment,
and \texttt{db/seed}, containing sample-data migrations that run only
in the \texttt{local} and \texttt{dev} Spring profiles. This split
guarantees that production never receives synthetic seed data. At the
time of writing, the schema has reached \textbf{V12}. Table
\ref{tab:migrations} summarizes the migration history.

\begin{longtable}{@{}p{0.20\linewidth}p{0.58\linewidth}p{0.12\linewidth}@{}}
\caption{Flyway Migration History (V1--V12)}
\label{tab:migrations}\\
\toprule
\textbf{Migration} & \textbf{Description} & \textbf{Environment} \\
\midrule
\endhead
V1\_\_initial\_schema & Base schema: users, posts, courses, course\_groups, group\_members, course\_posts, form\_templates, form\_submissions, course\_grades, stored\_files & All \\
V2\_\_admin\_users & admin\_users table, shared read access with myIU-admin & All \\
V3\_\_user\_provisioning & ALTER users: is\_active, ms\_tenant\_id, provisioned\_by, provisioned\_at & All \\
V4\_\_support\_tickets & support\_tickets table & All \\
V5\_\_login\_sessions & login\_sessions table (ip/country/city/browser/os/device\_type, is\_revoked) & All \\
V6\_\_course\_schedules & course\_schedules table & All \\
V7\_\_seed\_data & Baseline sample data (users, courses, groups) & local/dev \\
V8\_\_assignment\_submissions & assignment\_submissions table (course\_post\_id initially without a foreign key -- see Section 5.2) & All \\
V9\_\_attendance\_records & attendance\_records table & All \\
V10\_\_seed\_data\_extended & Extended sample data (group\_members, schedules, grades, attendance, submissions) & local/dev \\
V11\_\_cascade\_user\_deletes & 7 foreign keys changed from NO ACTION to CASCADE on users(id); adds the missing foreign key on assignment\_submissions.course\_post\_id & All \\
V12\_\_login\_session\_precise\_location & Adds latitude/longitude/province/district/ward to login\_sessions & All \\
\bottomrule
\end{longtable}

\section{Multi-Environment Configuration}

The system runs under three Spring profiles -- \texttt{local},
\texttt{dev}, and \texttt{prod} -- each adjusting a consistent set of
concerns: Flyway migration locations (whether \texttt{db/seed} is
included), the HikariCP connection pool size, the logging verbosity,
and the visibility of the Swagger UI, which is enabled on
\texttt{local}/\texttt{dev} for interactive API exploration and
disabled entirely on \texttt{prod}. All environment-specific and
sensitive values -- database credentials, the JWT signing secret, and
CORS allowed origins -- are injected through environment variables
rather than hardcoded in the repository, following the same convention
across both the Portal and Admin backends.

Figure \ref{fig:deployment} shows how the \texttt{prod} profile maps
onto the production Docker Compose topology: an Nginx reverse proxy
terminates TLS and routes by path to the four containers (Portal
frontend/backend, Admin frontend/backend), which all share the same
Docker bridge network and the single PostgreSQL container.

% SCREENSHOT SOURCE: docs/system-design.html, diagram 2.2
% "Deployment Diagram -- Docker Compose Production" (article id="s2-2")
\begin{figure}[htbp]
\centering
\includegraphics[width=0.9\linewidth]{diagram-2-2-deployment.png}
\caption{Deployment Diagram -- Docker Compose Production (source: \texttt{docs/system-design.html}, diagram 2.2)}
\label{fig:deployment}
\end{figure}

\chapter{Implementation and Results}

\section{Backend Implementation}

The entity layer uses Jakarta Persistence (JPA) with Lombok to minimize
boilerplate. All entities use \texttt{GenerationType.UUID} so that
identifier generation is controlled by the application rather than a
database trigger. The \texttt{User} entity illustrates this, along with
the pattern used throughout the schema for tracking how an account was
created:

\begin{quote}
\ttfamily\footnotesize
\noindent @Entity\\
@Table(name = "users")\\
public class User \{\\
\hphantom{xx}@Id @GeneratedValue(strategy = GenerationType.UUID)\\
\hphantom{xx}private UUID id;\\[2pt]
\hphantom{xx}@Column(unique = true, nullable = false)\\
\hphantom{xx}private String email;\\[2pt]
\hphantom{xx}private String password; // null for Microsoft SSO accounts\\
\hphantom{xx}private String authProvider = "local"; // "local" | "microsoft"\\
\hphantom{xx}private String msTenantId;\\[2pt]
\hphantom{xx}@ElementCollection(fetch = FetchType.EAGER)\\
\hphantom{xx}@CollectionTable(name = "user\_roles", joinColumns = @JoinColumn(name = "user\_id"))\\
\hphantom{xx}private Set<String> roles = new HashSet<>(); // "student" | "lecturer"\\
\}
\end{quote}

The \texttt{roles} collection is deliberately fetched eagerly, since
Spring Security must know a principal's roles at the moment of
authentication, before any lazy-loading opportunity would otherwise
arise; roles are stored as plain strings and prefixed to
\texttt{ROLE\_student} / \texttt{ROLE\_lecturer} inside
\texttt{UserDetails}, matching Spring Security's
\texttt{@PreAuthorize("hasRole(...)")} convention.

The service layer is where \texttt{@Transactional} boundaries matter
most. With \texttt{spring.jpa.open-in-view} disabled, the Hibernate
session closes as soon as the surrounding transaction ends, so any
method that later touches a \texttt{FetchType.LAZY} relationship must
itself be wrapped in a transaction, or Hibernate will throw a
\texttt{LazyInitializationException} once the session is gone [5].
\texttt{TimetableService.getForUser()} is a representative example: it
must access \texttt{schedule.getCourse()} inside its DTO-mapping step,
so the whole method is annotated \texttt{@Transactional(readOnly =
true)}, which both keeps the session open long enough for the lazy
proxy to resolve and signals to PostgreSQL that the transaction is
read-only.

Controllers are kept intentionally thin. \texttt{GroupMemberController}
supports five different combinations of optional query parameters
(\texttt{groupId}, \texttt{courseId}, \texttt{studentId}) through an
explicit if/else chain rather than a more "clever" generic query
builder -- a deliberate readability-over-generality choice, since every
branch corresponds to a distinct, predictable frontend use case.
\texttt{CourseGroupController.getLecturerCourses()} is a case where
\texttt{@Transactional} is placed directly on the controller method
rather than pushed down into a service: the endpoint is a small,
single-purpose two-line LAZY-relationship traversal that is not reused
elsewhere, and introducing a service method purely to wrap it would add
indirection without benefit.

\section{Frontend Implementation (Portal)}

The Portal frontend is built with React 18, Vite, and React Query v5
(\texttt{@tanstack/react-query}), which owns all server state so that
API data is never duplicated into local component state. A
representative hook:

\begin{quote}
\ttfamily\footnotesize
\noindent export const useGetTimetable = () =>\\
\hphantom{xx}useQuery(\{\\
\hphantom{xxxx}queryKey: ['timetable'],\\
\hphantom{xxxx}queryFn: getTimetable,\\
\hphantom{xxxx}staleTime: 5 * 60 * 1000, // 5 min, no refetch on tab focus\\
\hphantom{xx}\});
\end{quote}

\texttt{TimetablePage} renders the weekly schedule as a CSS grid with
six day columns and a fixed 64px-per-hour vertical scale
(\texttt{HOUR\_PX = 64}); a small \texttt{timeToPx()} helper converts a
lesson's start/end time directly into pixel offsets for absolute
positioning inside the grid. \texttt{AuthContext.tsx} manages the
authentication lifecycle end to end: it stores the JWT in
\texttt{localStorage}, injects it into every request through an Axios
interceptor, and, on sign-out, additionally redirects the browser to
Microsoft's own logout endpoint so that the Azure AD session itself is
also terminated, not just the local application state.

\section{Admin Panel Implementation}

The Admin backend lives in a separate repository, \texttt{myIU-admin},
deployed independently on port 8081 and connected to the same
PostgreSQL database with \texttt{ddl-auto: none}. As described in
Section 3.2.1, its early Excel-based bulk provisioning feature
(\texttt{POST /api/users/provision}, backed by Apache POI) has since
been removed entirely; administrators now create or deactivate accounts
one at a time through the Users screen, with each action still recorded
through the \texttt{provisioned\_by}/\texttt{provisioned\_at} audit
columns introduced in Section 3.3.

Because the Admin Panel is the system's highest-privilege surface, it
applies a full set of OWASP-recommended HTTP security headers [8]: a
Content Security Policy restricting script and style sources to
\texttt{'self'}, HTTP Strict Transport Security with a one-year max-age,
\texttt{frame-ancestors 'none'} paired with \texttt{X-Frame-Options:
DENY} to prevent clickjacking, a strict referrer policy, and a
Permissions Policy that explicitly revokes camera, microphone, and
geolocation access that the Admin frontend never needs.

\section{API Documentation and Testing}

Interactive API documentation is served through Swagger UI
(\texttt{/swagger-ui.html}) and Springdoc OpenAPI on the \texttt{local}
and \texttt{dev} profiles, and is disabled entirely on \texttt{prod}.
Rather than reporting testing in the abstract, Table
\ref{tab:test-results} reports the exact result of running each
project's own test and lint commands, on the date this thesis was
written -- not an estimate or a description of intended coverage.

\begin{longtable}{@{}p{0.28\linewidth}p{0.44\linewidth}p{0.20\linewidth}@{}}
\caption{Actual Test and Lint Results at Time of Writing}
\label{tab:test-results}\\
\toprule
\textbf{Project} & \textbf{Command} & \textbf{Result} \\
\midrule
\endhead
Portal backend (\texttt{myIU/backend}) & \texttt{./mvnw test} & 8 tests, 100\% pass (after fixing 1 stale test -- see Section 5.2) \\
Admin backend (\texttt{myIU-admin/backend}) & \texttt{./mvnw test} & No test directory exists -- 0 tests \\
Portal frontend (\texttt{myIU/frontend}) & \texttt{npm run lint} & 177 problems (165 errors, 12 warnings); project's own \texttt{--max-warnings 0} config means this build fails its own lint gate \\
Admin frontend (\texttt{myIU-admin/frontend}) & \texttt{npm run lint} & No lint script configured at all \\
\bottomrule
\end{longtable}

The Portal's continuous integration workflow (\texttt{.github/workflows/ci.yml})
runs \texttt{npx tsc --noEmit} plus lint on the frontend job and
\texttt{./mvnw test} on the backend job for every push and pull request
into \texttt{main}/\texttt{develop} -- so CI is correctly configured to
catch exactly the two problems above. What could not be confirmed
directly (no access to the GitHub Actions dashboard while preparing
this thesis) is whether CI had actually been green on \texttt{main}
before the most recent merges; running the same commands locally, at
the same point in time, reproduced red results for both jobs.

\section{Results}

At the time of writing, the following flows have been manually
exercised end to end and confirmed working: Microsoft SSO sign-in and
JWT issuance for both student and lecturer accounts; course, group, and
membership management; the weekly timetable grid; assignment submission
and grading with real-time WebSocket notification delivery; single and
bulk attendance marking; the component-based grade view; the academic
social feed (posts, comments, likes, saves); dynamic form submission
and administrator approval/rejection with email notification; support
ticket filing and resolution; and login-session listing, revocation,
and the opt-in GPS-based precise-location flow. No UI screenshots are
reproduced in this thesis, as they were not part of the source material
this rewrite was built from; they should be added from the running
application before final submission.

\chapter{Discussion and Evaluation}

\section{Discussion}

\textbf{OSIV disabled versus explicit \texttt{@Transactional}.} During
development, the timetable page repeatedly returned an empty list even
though the database clearly held matching rows. The root cause was the
interaction between \texttt{open-in-view: false} and a missing
\texttt{@Transactional} annotation on \texttt{TimetableService}: the
Hibernate session closed before the DTO-mapping step accessed
\texttt{schedule.getCourse()}, throwing a
\texttt{LazyInitializationException} that Spring Boot converted into an
HTTP 500, which the frontend silently rendered as an empty list.
Spring Boot's default (\texttt{open-in-view: true}) would have hidden
this problem by keeping the session open for the entire HTTP request --
a performance anti-pattern, since the database session then stays open
through JSON serialization -- but \texttt{open-in-view: false} is the
better default only if every developer consistently remembers exactly
where \texttt{@Transactional} is required, a convention that is easy to
violate silently and hard to catch without direct testing against a
real database.

\textbf{Dual-backend trade-off.} Splitting Portal and Admin into two
backends sharing one database buys separation of concerns, a smaller
attack surface for the higher-traffic Portal, and independent scaling,
at the cost of a schema-integrity boundary that is currently enforced
only by convention (\texttt{ddl-auto: none}) rather than by database
permissions -- nothing at the database level actually prevents the
Admin backend from issuing raw DDL if a developer ever added code that
did so.

\textbf{Stateless JWT versus server-side sessions.} A stateless JWT
avoids the operational cost of a shared session store and scales
horizontally without sticky sessions, but it requires an explicit
revocation mechanism, addressed here with a two-layer design (in-memory
blacklist plus a persistent \texttt{login\_sessions} table) rather than
relying on token expiry alone.

\textbf{Removing the Excel provisioning feature.} As discussed in
Section 3.2.1, the decision to remove Excel-based bulk provisioning
entirely, rather than only deleting its unreachable Portal-side dead
code, weighed the attack surface and dependency cost of a rarely used
feature against the real operational cost of provisioning accounts one
at a time -- judged acceptable at the system's current scale.

\textbf{Virtual threads.} The Portal enables Project Loom virtual
threads (\texttt{spring.threads.virtual.enabled: true}) [9]. For an
I/O-bound workload typical of a university portal -- many small
requests, most of the latency spent waiting on the database or an
external API -- virtual threads let each request occupy a lightweight,
JVM-managed thread rather than a dedicated operating-system thread,
which should improve throughput under concurrent load without a
corresponding increase in OS thread count, though no formal benchmark
of this effect has been run (see Section 5.2.2).

\section{Evaluation}

\subsection{Security and Reliability}

Table \ref{tab:security-findings} lists concrete defects found and
fixed during development, each verified by reproducing the problem
directly rather than assumed from a code read alone.

\begin{longtable}{@{}p{0.05\linewidth}p{0.30\linewidth}p{0.32\linewidth}p{0.23\linewidth}@{}}
\caption{Security and Reliability Findings}
\label{tab:security-findings}\\
\toprule
\textbf{\#} & \textbf{Issue} & \textbf{How It Was Found} & \textbf{Status} \\
\midrule
\endhead
1 & IP spoofing via \texttt{X-Forwarded-For} -- bypasses login rate-limit and forges the audit-log location & Manual comparison of the header value against the displayed location & Fixed; trust-boundary check added (Section 3.6.1) with a regression test \\
2 & \texttt{@Async} + \texttt{@Modifying} query missing \texttt{@Transactional} -- 100\% of GeoIP enrichment silently failed & Direct database query showed every session stuck at \texttt{country = "Resolving"}; confirmed via \texttt{TransactionRequiredException} in logs & Fixed \\
3 & Case-sensitive email comparison caused legitimate SSO logins to fail (\texttt{not\_provisioned}) & Reproduced the real login failure; compared stored email against the OAuth claim & Fixed (data correction, 4 rows) \\
4 & Missing foreign key on \texttt{assignment\_submissions.course\_post\_id} -- orphaned rows possible & Discovered while testing cascade-delete end to end & Fixed (migration V11) \\
5 & Portal and Admin used two separate \texttt{uploads/} directories while generated URLs assumed a shared one & User-reported "download does nothing"; inspection found the Admin directory empty & Fixed (single absolute, shared \texttt{UPLOAD\_DIR}) \\
6 & \texttt{String.format("\%f", ...)} is locale-sensitive and can corrupt the Nominatim query URL on a comma-decimal JVM locale & Proactive code review after feature completion & Fixed (\texttt{Locale.US} pinned) \\
7 & Duplicated ward/district text in reverse-geocoded addresses & Direct user report & Fixed (deduplication, Section 3.7) \\
\bottomrule
\end{longtable}

Every one of these seven issues was found through manual
investigation \emph{after} a symptom was observed -- a failed login, an
empty result set, a user complaint -- rather than by an automated test
or static analysis tool catching it beforehand. That pattern is itself
the clearest evidence for the testing gap quantified in the next
subsection.

\subsection{Testing and Code Quality}

Section 4.4 already reported the concrete numbers; this subsection
interprets them. The Portal backend's entire automated test suite
consists of two classes (\texttt{GeoIpServiceTest},
\texttt{SessionServiceTest}), both unit tests using Mockito, with no
integration test that touches a real database -- despite CI
documentation referencing Testcontainers, no test in the repository
currently uses it. One of those eight tests
(\texttt{createSession\_uses\_x\_forwarded\_for\_header}) was found
\emph{failing} the first time it was run for this evaluation: it had
been written before the IP-spoofing fix in Table
\ref{tab:security-findings} and was never updated afterward, meaning
the suite had been red since that security fix landed without anyone
re-running it. It was corrected and a dedicated regression test for the
spoofing fix itself was added in the same pass. The Admin backend has
no test directory at all, which is the single highest-risk gap in the
system: its cascade-delete behavior, foreign-key-violation handling,
and session-revocation logic have no automated safety net whatsoever.

On the frontend, the Portal's own \texttt{package.json} configures
\texttt{eslint --max-warnings 0}, meaning that, by the project's own
definition of a passing build, its current 177 lint problems (165
errors, 12 warnings, mostly \texttt{@typescript-eslint/no-explicit-any}
and unused variables) constitute a failing build, not an acceptable
level of warnings. The Admin frontend has no lint script configured at
all, so its code quality has never been automatically checked.

\subsection{Strategic Technical Debt and Trade-offs}

Reported plainly, without smoothing over the gaps: there is
\textbf{no SAST or SCA scan} (no SonarQube, no Snyk) on either
repository, so no dependency-vulnerability data exists for either
codebase; there is \textbf{no performance or load benchmark}, so no
real latency or throughput figures can be cited for the virtual-thread
claim in Section 5.1; there is \textbf{no integration test that touches
a real database}; the \textbf{Admin backend has zero tests}; the
\textbf{Admin frontend has no lint configuration}; and \textbf{CI is
correctly written but not confirmed to have been green} on the most
recent merges into \texttt{main} during this evaluation. These points
are stated directly because they are precisely the details most easily
omitted when a report only describes which features work on the UI,
and a report that hides them would misrepresent the system's actual
state of readiness.

The Client IP Trust Boundary described in Section 3.6.1 also currently
infers "trusted" purely from the peer address being private or
loopback, rather than from an explicit configured allow-list -- adequate
for a single-host deployment but insufficient for a more complex
multi-proxy topology, as discussed further in Section 6.2. Likewise,
the GPS reverse-geocoding described in Section 3.7 is explicitly
best-effort: OpenStreetMap's Vietnamese address data is not uniformly
structured across all three administrative levels everywhere in the
country after recent mergers, and the system resolves that
inconsistency by omitting a level rather than guaranteeing all three are
always present.

\chapter{Conclusion and Future Work}

\section{Conclusion}

This thesis presented the design and implementation of myIU, a
university student portal that delegates authentication entirely to
Microsoft Azure AD while retaining full ownership of academic data
through a purpose-built PostgreSQL schema, and that separates
student/lecturer-facing functionality from administrative functionality
behind two independently deployed Spring Boot backends. The system
covers course and group management, timetabling, attendance,
assignment submission and grading, an academic social feed,
form-based administrative requests, and login-session monitoring with
an opt-in precise-location feature.

Beyond the working feature set, this thesis makes the following
technical contributions: (1) a dual-backend architecture with a clearly
delineated schema-ownership boundary between a Flyway-driven Portal and
a schema-frozen Admin Panel; (2) a two-layer JWT revocation design
combining an in-memory blacklist with a persistent, database-backed
session table; (3) a Flyway multi-location migration strategy that
separates schema evolution from environment-specific seed data; (4) a
documented pattern for handling \texttt{LazyInitializationException}
correctly under \texttt{open-in-view: false}; (5) a deliberately
classified cascade-delete policy distinguishing self-owned data from
content that other users depend on; (6) a client IP trust-boundary
mechanism that closes a real spoofing vulnerability in both the
rate-limiter and the audit log; and (7) an opt-in GPS reverse-geocoding
flow, together with the discovery and fix of a subtle
\texttt{@Transactional} + \texttt{@Async} interaction that had silently
left every login session's location unresolved for an extended period
without producing a single visible error.

Chapter 5 additionally evaluated the system against its own test suite
and linters exactly as they stand today, rather than only against its
feature list, and reported concrete numbers -- eight passing backend
tests after fixing one that had gone stale, zero tests on the Admin
backend, and a frontend lint run that fails the project's own quality
gate -- because an honest account of a system's readiness has to
include what has not yet been verified, not only what has been built.

\section{Future Work and Technical Debt Mitigation}

The following priorities are identified for the next phase of
development, in the order suggested by the evaluation in Chapter 5:

\begin{itemize}
\item \textbf{Enforce CI as a merge gate.} Confirm the existing
\texttt{test-portal} job is green and make it a required branch
protection check, so that a change cannot be merged into \texttt{main}
while the backend test suite is red.
\item \textbf{Test the Admin backend}, starting with its two
highest-risk areas: cascade-delete behavior and foreign-key-violation
handling, both of which currently have no automated coverage at all.
\item \textbf{Run a free-tier SAST/SCA scan} (e.g.\ Snyk) across both
repositories -- a near-zero-cost action that would close an
information gap that currently exists for both codebases.
\item \textbf{Bring the Portal frontend back under its own lint gate}
by resolving the 177 outstanding ESLint problems down to the
\texttt{--max-warnings 0} threshold the project already declares for
itself, and configure a lint script for the Admin frontend, which
currently has none.
\item \textbf{Distributed JWT blacklist.} Replace the in-memory
\texttt{ConcurrentHashMap} blacklist with a shared store such as Redis,
so that token revocation is consistent across horizontally scaled
backend instances rather than local to a single node.
\item \textbf{Database-level schema protection.} Grant the Admin
backend a database role limited to DML (\texttt{SELECT},
\texttt{INSERT}, \texttt{UPDATE}, \texttt{DELETE}) without DDL
privileges, turning the current \texttt{ddl-auto: none} convention into
a boundary the database itself enforces.
\item \textbf{Explicit trusted-proxy allow-list.} Replace the current
private/loopback heuristic in the IP trust boundary (Section 3.6.1)
with an explicit, environment-variable-configured allow-list of proxy
addresses, needed for any deployment topology more complex than a
single host.
\item \textbf{Integration testing against a real database.} Adopt
Testcontainers PostgreSQL for at least the highest-risk service
methods -- particularly those relying on \texttt{@Transactional} and
LAZY loading -- so that defects such as the ones in Table
\ref{tab:security-findings} are caught before a symptom is observed in
practice, not after.
\item \textbf{Rebuild the deployment pipeline correctly.} The existing
deployment workflows assumed \texttt{myIU-admin} was a subdirectory of
this repository, when it is in fact a separate repository; they have
never successfully deployed and should be redesigned for the correct
two-repository topology (either a dual checkout in one pipeline or a
separate pipeline per repository) rather than left in a broken state.
\item \textbf{Performance benchmarking and observability.} Establish a
real load-testing baseline to substantiate the virtual-threads
performance discussion in Section 5.1, and build out the
\texttt{/actuator/prometheus} endpoint already exposed by the backend
into a proper Grafana dashboard.
\item \textbf{Mobile application.} A React Native client reusing the
existing Portal REST API would extend the platform to native mobile
without backend changes.
\end{itemize}

% ------------------------------------------------------------------ %
% REFERENCES
% ------------------------------------------------------------------ %
\chapter*{References}
\addcontentsline{toc}{chapter}{References}

\begin{enumerate}
\item Hardt, D. (2012). \emph{The OAuth 2.0 Authorization Framework}. RFC 6749. IETF.
\item Sakimura, N., Bradley, J., Jones, M., de Medeiros, B., \& Mortimore, C. (2014). \emph{OpenID Connect Core 1.0}. OpenID Foundation.
\item Jones, M., Bradley, J., \& Sakimura, N. (2015). \emph{JSON Web Token (JWT)}. RFC 7519. IETF.
\item Spring Framework Team. (2024). \emph{Spring Boot Reference Documentation 3.4.x}. VMware Tanzu / Broadcom.
\item Hibernate ORM Team. (2024). \emph{Hibernate ORM 6.6 User Guide: Fetching Strategies}. Red Hat.
\item Moodle. (2024). \emph{Moodle Statistics}. Retrieved from moodle.org.
\item Instructure. (2024). \emph{Canvas LMS Developer Documentation}. Instructure Inc.
\item OWASP Foundation. (2024). \emph{HTTP Security Response Headers Cheat Sheet}. OWASP.
\item Lowe, D. (2023). \emph{Project Loom: Fibers and Continuations for the Java Virtual Machine}. OpenJDK.
\item Flyway Team. (2024). \emph{Flyway Documentation: Concepts and Configuration}. Redgate.
\item Anthology Inc. (2024). \emph{Blackboard Learn Administrator and Developer Documentation}. Anthology Inc.
\item Microsoft. (2024). \emph{Microsoft Teams for Education Documentation}. Microsoft Learn.
\end{enumerate}

\end{document}
